\chapter{2011年上海卷（春）}

\section{填空题}

\begin{problem}
函数 \( y = \lg (x - 2) \) 的定义域是 \fillinblank{} \fillinblank{}.
\end{problem}

\begin{problem}
若集合 \(A = \{x \mid x \geqslant 1\}\)，\(B = \{x \mid x^2 \leqslant 4\}\)，则 \(A \cap B = \)  \fillinblank{}.
\end{problem}

\begin{problem}
在 \(\triangle ABC\) 中，若 \(\tan A = \frac{\sqrt{2}}{3}\)，则 \(\sin A = \)  \fillinblank{}.
\end{problem}

\begin{problem}
若行列式 \( \begin{vmatrix}2^{x}&4\\1&2\end{vmatrix}=0 \) ，则 x=\fillinblank{} \fillinblank{}.
\end{problem}

\begin{problem}
若 \( \sin x = \frac{1}{3}, x \in \left[-\frac{\pi}{2}, \frac{\pi}{2}\right] \) ，则 \( x = \) （结果用反三角函数值表示） \fillinblank{}.
\end{problem}

\begin{problem}
\(\left(x + \frac{1}{x}\right)^6\) 的二项展开式的常数项为 \fillinblank{} \fillinblank{}.
\end{problem}

\begin{problem}
两条直线 \(l: x - \sqrt{3} y + 2 = 0\) 与 \(l_2: x - y + 2 = 0\) 夹角的大小是 \fillinblank{} \fillinblank{}.
\end{problem}

\begin{problem}
若 \( S_{n} \) 为等比数列 \( \{a_{n}\} \) 的前 n 项和，\( 8a_{2} + a_{5} = 0 \) ，则 \( \frac{8a}{S_{3}} = \)  \fillinblank{}.
\end{problem}

\begin{problem}
若椭圆 C 的焦点和顶点分别是双曲线 \( \frac{x^{2}}{5}-\frac{y^{2}}{4}=1 \) 的顶点和焦点，则椭圆 C 的方程是 \fillinblank{} \fillinblank{}.
\end{problem}

\begin{problem}
若点 \(O\) 和点 \(F\) 分别为椭圆 \(\frac{x^2}{2} + y^2 = 1\) 的中心和左焦点，点 \(P\) 为椭圆上的任意一点，则 \(|OP|^2 + |PF|^2\) 的最小值为 \fillinblank{} \fillinblank{}.
\end{problem}

\begin{problem}
根据如图所示的程序框图，输出结果 i =\fillinblank{} \fillinblank{}.
\end{problem}

\begin{problem}
2011年上海春季高考有8所高校招生，如果某3位同学恰好被其中2所高校录取，那么录取方法的种数为 \fillinblank{} \fillinblank{}.
\end{problem}

\begin{problem}
有一种多面体的饰品，其表面由6个正方形和8个正三角形组成(如图), \(AB\) 与 \(CD\) 所成角的大小是 \fillinblank{}.
\end{problem}

\begin{problem}
为求解方程 \(x^{5} - 1 = 0\) 的虚根，可以把原方程变形为 \((x - 1)(x^4 + x^3 + x^2 + x + 1) = 0\)，再变形为 \((x - 1)(x^2 + ax + 1)(x^2 + bx + 1) = 0\)，由此可得原方程的一个虚根为 \fillinblank{} \fillinblank{}.
\end{problem}

\section{单选题}

\begin{problem}
若向量 \(\overrightarrow{a} = (2,0)\)，\(\overrightarrow{b} = (1,1)\)，则下列结论正确的是（\quad）

\choices
    {\(\overrightarrow{a} \cdot \overrightarrow{b} = 1\).}
    {\(|\overrightarrow{a}| = |\overrightarrow{b}|\).}
    {\((\overrightarrow{a} -\overrightarrow{b})\bot \overrightarrow{b}\).}
    {\(\overrightarrow{a} // \overrightarrow{b}\).}

\end{problem}

\begin{answer}
C
\end{answer}
\begin{problem}
函数 \( f(x) = \frac{4x - 1}{2^x} \) 的图像关于（\quad）

\choices
    {原点对称.}
    {直线 y = x 对称.}
    {直线 y = -x 对称.}
    {y 轴对称.}

\end{problem}

\begin{answer}
A
\end{answer}
\begin{problem}
直线 \(l: y = k\left(x + \frac{1}{2}\right)\) 与圆 \(C: x^2 + y^2 = 1\) 的位置关系为（\quad）

\choices
    {相交或相切.}
    {相交或相离.}
    {相切.}
    {相交.}

\end{problem}

\begin{answer}
D
\end{answer}
\begin{problem}
若 \(\overrightarrow{a_1}, \overrightarrow{a_2}, \overrightarrow{a_3}\) 均为单位向量，则 \(\overrightarrow{a_1} = \left(\frac{\sqrt{3}}{3}, \frac{\sqrt{6}}{3}\right)\) 是 \(\overrightarrow{a_1} + \overrightarrow{a_2} + \overrightarrow{a_3} = (\sqrt{3}, \sqrt{6})\) 的（\quad）

\choices
    {充分不必要条件.}
    {必要不充分条件.}
    {充分必要条件.}
    {既不充分又不必要条件.}

\end{problem}

\begin{answer}
B
\end{answer}
\section{解答题}

\begin{problem}
已知向量 \( \overrightarrow{a} = (\sin 2x - 1, \cos x) \) ， \( \overrightarrow{b} = (1, 2\cos x) \) . 设函数 \( f(x) = \overrightarrow{a} \cdot \overrightarrow{b} \) ，求函数 \( f(x) \) 的最小正周期及 \( x \in \left(0, \frac{\pi}{2}\right] \) 时的最大值. 
\end{problem}

\begin{problem}
某甜品店制作一种蛋筒冰淇淋，其上半部分呈半球形、下半部分呈圆锥形(如图). 现把半径为 10cm 的圆形蛋皮等分成 5 个扇形，用一个扇形蛋皮围成圆锥的侧面(蛋皮厚度忽略不计)，求该蛋筒冰淇淋的表面积和体积(精确到 0.01).
\end{problem}

\begin{problem}
已知抛物线 \( F: x^{2} = 4y \) . (1) \(\triangle ABC\) 的三个顶点在抛物线 \(F\) 上，记: \(\triangle ABC\) 的三边 \(AB, BC, CA\) 所在直线的斜率分别为 \(k_{AB}, k_{BC}, k_{CA}\)，若点 \(A\) 在坐标原点，求 \(k_{AB} - k_{BC} + k_{CA}\) 的值; (2) 请你给出一个以 \(P(2,1)\) 为顶点、且其余各顶点均为抛物线 \(F\) 上的动点的多边形，写出多边形各边所在直线的斜率之间的关系式，并说明理由.
\end{problem}

\begin{problem}
定义域为 \(\R\)，且对任意实数 \(x_{1}, x_{2}\) 都满足不等式 \(f\left(\frac{x_{1} + x_{2}}{2}\right) \leqslant \frac{f(x_{1}) + f(x_{2})}{2}\) 的所有函数 \(f(x)\) 组成的集合记为 \(M\). 例如，函数 \(f(x) = kx + b \in M\). (1) 已知函数 \( f(x) = \begin{cases} x, & x \geqslant 0, \\ \frac{1}{2} x, & x < 0. \end{cases} \) 证明：\( f(x) \in M; \) (2) 写出一个函数 \( f(x) \)，使得 \( f(x) \notin M \)，并说明理由； (3) 写出一个函数 \( f(x) \in M \)，使得数列极限 \( \lim_{n \to \infty} \frac{f(n)}{n^2} = 1, \lim_{n \to \infty} \frac{f(-n)}{-n} = 1. \)
\end{problem}

\begin{problem}
对于给定首项 \(x_0 > \sqrt[3]{a}\)（\(a > 0\)），由递推式 \(x_{n+1} = \frac{1}{2}\left(x_n + \sqrt{\frac{a}{x_n}}\right)\)（\(n \in \mathbb{N}\)）得到数列 \(\{x_n\}\)，且对于任意的 \(n \in \mathbb{N}\)，都有 \(x_n > \sqrt[3]{a}\). 用数列 \(\{x_n\}\) 可以计算 \(\sqrt[3]{a}\) 的近似值. （1）取 \(x_0 = 5, a = 100\)，计算 \(x_1, x_2, x_3\) 的值（精确到 0.01）；归纳出 \(x_n, x_{n+1}\) 的大小关系；（2）当 \(n \geqslant 1\) 时，证明 \(x_n - x_{n+1} < \frac{1}{2}(x_{n-1} - x_n)\)；（3）当 \(x_0 \in [5, 10]\) 时，用数列 \(\{x_n\}\) 计算 \(\sqrt[3]{100}\) 的近似值，要求 \(|x_n - x_{n+1}| < 10^{-4}\)，请你估计 \(n\)，并说明理由.
\end{problem}

